\documentclass[12pt]{amsart}
\usepackage[margin=1in]{geometry}
\usepackage{amsmath,amssymb,amsthm}
\usepackage{mathtools}
\usepackage{hyperref}

\theoremstyle{plain}
\newtheorem{theorem}{Theorem}[section]
\newtheorem{proposition}[theorem]{Proposition}
\newtheorem{corollary}[theorem]{Corollary}
\newtheorem{lemma}[theorem]{Lemma}
\newtheorem{definition}[theorem]{Definition}
\theoremstyle{remark}
\newtheorem*{remark}{Remark}

\title[The Orbital Phase Equation on the AM--GM Cone]{Closing the Clock:\\
An Exact Orbital-Phase Equation on the AM--GM Cone}
\author{Jeremy Ebert}
\date{2026 (v2)}

\begin{document}
\maketitle

\begin{abstract}
The AM--GM cone construction of \cite{Ebert2026Cone} represents an orbit's turning points $(r_a,r_p)$ as a point $(X,Y,T)$ on the quadric $X^2+Y^2=T^2$, and shows that total energy, angular momentum, and central mass isolate along coordinate directions of that cone. That construction is a description of orbital \emph{shape} only: it fixes $(a,e)$, equivalently $(X,Y,T)$, but says nothing about where the body sits on its orbit at a given time. We supply the missing piece. Writing the orbital radius $r$ directly in cone coordinates as $r=Y^2/(T+X\cos\nu)$, and combining this with the cone's own angular-momentum formula $h=\sqrt{\mu Y^2/T}$, Kepler's equal-area law $r^2\dot\nu=h$ becomes an explicit, closed-form ODE for the true anomaly $\nu(t)$ in terms of $(X,Y,T)$ alone. This equation couples directly to the shape-evolution equations already derived for perturbed motion in \cite{Ebert2026Cone}, giving a single coupled system $(\dot X,\dot Y,\dot T,\dot\nu)$ that tracks both the orbit's shape and its instantaneous position, closing the gap explicitly left open there. We verify the unperturbed case numerically against direct Cartesian integration of the two-body equations of motion, with agreement to floating-point precision. We then connect the clock to the algebraic machinery of the companion papers \cite{Ebert2026Table,Ebert2026EigenCoord} in two ways: we show the classical eccentric-to-true-anomaly map is exactly a hyperbolic Möbius transformation of the phase circle, at rapidity half the orbit's own $\operatorname{artanh}(e)$, generated by the cone's own $B_X$; and we show the shape-flow equations $(\dot X,\dot Y,\dot T)$ collapse to a single split-quaternion commutator equation, while proving algebraically that the clock variable cannot join that same equation --- an exact, structural confirmation of the gap this note exists to fill.
\end{abstract}

\section{The Gap}

The dynamic slicing operator of \cite{Ebert2026Cone} evolves the point $(X,Y,T)\in\mathbb{R}^3$ under a Lie-algebra flow in $\mathfrak{so}(2,1)$, and this is exactly the orbit's osculating \emph{shape}: $T=a$, $X=ae$, $Y=b$. Nothing in that construction determines the body's position on the ellipse at a given time $t$; $(X,Y,T)$ is silent on the true anomaly $\nu$, the classical clock variable of the two-body problem. This note adds that variable, expressed in the same coordinates, so that the cone construction and ordinary orbital timekeeping sit in one coupled system rather than two disconnected ones.

\section{The Radius in Cone Coordinates}

Recall the standard orbit equation $r=\dfrac{a(1-e^2)}{1+e\cos\nu}$, where $\nu$ is the true anomaly measured from periapsis.

\begin{proposition}[Radius as a function of $(X,Y,T,\nu)$]
Under $a=T$, $e=X/T$, and using the cone identity $T^2-X^2=Y^2$,
\begin{equation}
r(X,Y,T,\nu)=\frac{Y^2}{T+X\cos\nu}. \tag{14}
\end{equation}
In particular $r(\nu=0)=T-X=r_p$ and $r(\nu=\pi)=T+X=r_a$, recovering the two turning points the cone was built from in the first place.
\end{proposition}

\begin{proof}
Substituting $a=T$, $e=X/T$ into the orbit equation gives
\[
r=\frac{T\left(1-\frac{X^2}{T^2}\right)}{1+\frac{X}{T}\cos\nu}=\frac{T^2-X^2}{T+X\cos\nu}.
\]
Since $T^2-X^2=Y^2$ on the cone, this is exactly (14). At $\nu=0$: $r=(T^2-X^2)/(T+X)=T-X$. At $\nu=\pi$: $r=(T^2-X^2)/(T-X)=T+X$.
\end{proof}

\section{The Orbital Phase Equation}

\begin{theorem}[The clock equation]
Let $(X(t),Y(t),T(t))$ be any state on the cone $X^2+Y^2=T^2$ evolving under a (possibly time-dependent) flow, and let $\nu(t)$ denote the corresponding true anomaly. Then
\begin{equation}
\dot\nu \;=\; \frac{\sqrt{\mu}\,\bigl(T+X\cos\nu\bigr)^{2}}{Y^{3}\sqrt{T}}. \tag{15}
\end{equation}
\end{theorem}

\begin{proof}
Kepler's second law states $r^2\dot\nu=h$, with $h$ the (possibly time-dependent, if the orbit is perturbed) specific angular momentum. By Proposition 3.1(b) of \cite{Ebert2026Cone}, $h=\sqrt{\mu Y^2/T}$. Substituting this together with (14):
\[
\dot\nu=\frac{h}{r^2}=\frac{\sqrt{\mu Y^2/T}}{\left(\dfrac{Y^2}{T+X\cos\nu}\right)^{2}}=\frac{\sqrt{\mu}\,Y\,(T+X\cos\nu)^2}{\sqrt{T}\,Y^4}=\frac{\sqrt{\mu}\,(T+X\cos\nu)^2}{Y^3\sqrt{T}}. \qedhere
\]
\end{proof}

\begin{corollary}[The unperturbed case]
If $(X,Y,T)$ are constant (the unperturbed Kepler problem, Section 5 of \cite{Ebert2026Cone}), (15) is a single first-order ODE in $\nu$ alone, and integrating it reproduces ordinary Keplerian motion exactly: it is Kepler's equal-area law written in cone coordinates rather than in $(a,e)$, not a new dynamical law.
\end{corollary}

\begin{corollary}[The perturbed, coupled system]
Under a perturbing acceleration realized as a flow $\dot X, \dot Y, \dot T$ via equations (8)--(13) of \cite{Ebert2026Cone} (e.g.\ the atmospheric-drag generator coefficients (13)), equation (15) couples directly to that flow through $r=r(X,Y,T,\nu)$: the full state $(X,Y,T,\nu)$, evolving under
\[
\dot X = \text{(as in \cite{Ebert2026Cone})},\quad
\dot Y = \text{(as in \cite{Ebert2026Cone})},\quad
\dot T = \text{(as in \cite{Ebert2026Cone})},\quad
\dot\nu = \frac{\sqrt{\mu}\,(T+X\cos\nu)^2}{Y^3\sqrt T},
\]
is a closed, self-contained system that tracks both the osculating shape of the orbit and its instantaneous position along that orbit. This is precisely the missing half of the dynamic slicing operator: (14)--(15) supply the clock that \cite{Ebert2026Cone} leaves external to the construction.
\end{corollary}

\begin{remark}[What this is, and is not]
Equation (15) is algebraically equivalent to the standard true-anomaly rate equation of two-body mechanics, $\dot\nu=h/r^2$ with $h,r$ expressed in the usual orbital elements; nothing here is new orbital mechanics. What is supplied is the explicit rewriting of that equation entirely in terms of $(X,Y,T)$, so that it can be integrated alongside the cone's own shape-evolution equations rather than requiring a separate $(a,e)\to\nu$ bookkeeping step outside the construction. We are not aware of a stated closed-form clock equation for the AM--GM cone construction prior to this note.
\end{remark}

\section{Numerical Verification}

We integrate (15) for a fixed, unperturbed orbit ($\mu=398600.4418\ \mathrm{km^3/s^2}$, $a=24482\ \mathrm{km}$, $e=0.72$, starting at periapsis) and compare the resulting $\nu(t)$ against direct numerical integration of the Cartesian two-body equations of motion $\ddot{\mathbf r}=-\mu\mathbf r/r^3$, an independent ground truth not derived from (14)--(15).

\begin{center}
\begin{tabular}{lcc}
\hline
Quantity & Value \\
\hline
$\nu(t{=}20{,}000\ \mathrm{s})$ from (15) & $3.17792779558008$ rad \\
$\nu(t{=}20{,}000\ \mathrm{s})$ from Cartesian integration & $3.17792779558011$ rad \\
Difference & $-2.9\times10^{-14}$ rad \\
\hline
\end{tabular}
\end{center}

Agreement is at the level of floating-point roundoff, confirming (14)--(15) exactly reproduce two-body motion in the unperturbed case.

\section{The Eccentric Anomaly as a Hyperbolic Phase Boost}

Sections 2--4 give a closed clock equation for the true anomaly $\nu$, but $\nu$ is not the only classical clock variable in use: the eccentric anomaly $E$, defined by $r=T(1-e\cos E)$, is the variable in which Kepler's equation $M=E-e\sin E$ is linear in the mean anomaly $M=\sqrt{\mu/T^3}\,(t-t_0)$. This section records an exact identity connecting $E$ and $\nu$ using only $(X,Y,T)$, and shows that this identity is itself an action of the cone's own symmetry group \cite{Ebert2026EigenCoord} on a different space --- the phase circle rather than the cone --- giving a second point of contact between orbital timekeeping and the $\mathfrak{so}(2,1)$ structure underlying this whole project.

\begin{proposition}[Eccentric anomaly as a phase-circle boost]\label{prop:Eboost}
Let $(X,Y,T)$ be a point on the cone $X^2+Y^2=T^2$, $e=X/T$ the corresponding eccentricity, and $E,\nu$ the eccentric and true anomalies of the orbit. Then
\begin{equation}
\tan(\nu/2) = k\,\tan(E/2), \qquad k=\sqrt{\frac{T+X}{T-X}}=e^{\operatorname{artanh}(X/T)}. \tag{16}
\end{equation}
Writing $z=\tan(\theta/2)$ for the standard coordinate on the phase circle $\mathbb{RP}^1$, (16) is the action $z\mapsto kz$ of the $SL(2,\mathbb R)$ matrix
\begin{equation}
D(s)=\operatorname{diag}\!\big(e^{s/2},e^{-s/2}\big), \qquad s=\operatorname{artanh}(e), \tag{17}
\end{equation}
a hyperbolic Möbius transformation with rapidity exactly half the orbit's own $\operatorname{artanh}(e)$. Its two fixed points, $z=0$ and $z=\infty$, are exactly $\nu=E=0$ (periapsis) and $\nu=E=\pi$ (apoapsis).
\end{proposition}

\begin{proof}
Substituting $e=X/T$ into the classical relation $\tan(\nu/2)=\sqrt{(1+e)/(1-e)}\,\tan(E/2)$ turns the ratio under the root into $(T+X)/(T-X)$, giving (16). A diagonal matrix $\operatorname{diag}(a,a^{-1})\in SL(2,\mathbb R)$ acts on the coordinate $z$ of $\mathbb{RP}^1$ by the Möbius formula $z\mapsto (az)/a^{-1}=a^2z$; matching $a^2=k=e^s$ forces $a=e^{s/2}$, giving exactly (17), and then $z\mapsto a^2z=e^sz=kz$ reproduces (16). Diagonal matrices fix $z=0$ and $z=\infty$ identically; since $z=\tan(\theta/2)=0$ at $\theta=0$ and $z\to\infty$ at $\theta=\pi$, these are exactly $\nu=E=0$ and $\nu=E=\pi$, matching the turning points $r_p=T-X$, $r_a=T+X$ of Proposition 2.1. \qedhere
\end{proof}

\begin{remark}[Relation to the cone's own symmetry algebra]
The rapidity $s=\operatorname{artanh}(e)$ in (17) is exactly the parameter needed to reach eccentricity $e$ from the circular orbit $e=0$ under the $B_X$-flow of \cite{Ebert2026EigenCoord}: started from $p_0=q_0$, the flow $p(\phi)=p_0e^\phi,\,q(\phi)=q_0e^{-\phi}$ reaches $p/q=e^{2\phi}=(1+e)/(1-e)$ exactly at $\phi=\operatorname{artanh}(e)=s$. The Möbius map (17) uses \emph{half} this rapidity, $s/2$ --- the same halving that separates a $3$-dimensional vector representation of $\mathfrak{so}(2,1)$ from its $2$-dimensional $SL(2,\mathbb R)$ spin representation, and that already had to be imposed, for an unrelated reason, in the two-mode-squeezing correspondence of the companion quantization note. It is not a new coincidence; it is the same structural halving appearing here for a third time in this project.
\end{remark}

\subsection*{Numerical Verification}

We check (16) against the orbit and epoch already used in Section 4 ($\mu=398600.4418\ \mathrm{km^3/s^2}$, $a=T=24482$ km, $e=X/T=0.72$), using the value $\nu(t{=}20{,}000\,\mathrm s)=3.17792779558008$ rad computed there.

\begin{center}
\begin{tabular}{lc}
\hline
Quantity & Value \\
\hline
$E$ recovered from $\nu$ & $-3.05158767468174$ rad \\
$k=\sqrt{(T+X)/(T-X)}$ & $2.47847879612821$ \\
$\tan(\nu/2)$ & $-55.0370755471039075740743$ \\
$k\tan(E/2)$ & $-55.0370755471039075740743$ \\
Difference & $6.3\times10^{-29}$ \\
\hline
\end{tabular}
\end{center}

Agreement is at the level of the working precision (30 digits), confirming (16) holds exactly, not merely approximately, on a genuine two-body orbit.

\section{A Quaternion Form for the Shape Flow, and Why the Clock Cannot Join It}

Section 1 describes the cone's dynamics as ``a Lie-algebra flow in $\mathfrak{so}(2,1)$.'' This section makes that description completely literal: the three coupled scalar equations for $(\dot X,\dot Y,\dot T)$ collapse to a single split-quaternion commutator equation. We then show, algebraically rather than merely by the absence of a construction, that the clock variable cannot be absorbed into that same equation --- turning the gap identified in Section 1 into a proved structural fact.

\begin{definition}[Split-quaternion dictionary]
Let $\mathbb H'$ be the split-quaternion algebra with basis $1,i,j,k$, $i^2=-1$, $j^2=k^2=1$, $ij=k=-ji$, $jk=-i=-kj$, $ki=j=-ik$ (equivalently $ijk=1$). Identify the generators of the symmetry algebra of \cite{Ebert2026EigenCoord} with half the imaginary units of $\mathbb H'$:
\begin{equation}
L=\tfrac12 i,\qquad B_X=\tfrac12 j,\qquad B_Y=\tfrac12 k. \tag{18}
\end{equation}
\end{definition}

\begin{proposition}
Under (18), the commutator bracket of $\mathbb H'$ reproduces $[L,B_X]=B_Y$, $[L,B_Y]=-B_X$, $[B_X,B_Y]=-L$ exactly.
\end{proposition}

\begin{proof}
Direct computation from the multiplication table of $\mathbb H'$; verified symbolically, with zero residual in each bracket.
\end{proof}

\begin{theorem}[Shape flow as a quaternion commutator]\label{thm:laxshape}
Embed the state $(X,Y,T)$ as the pure split quaternion
\begin{equation}
V = T\,i + Y\,j - X\,k. \tag{19}
\end{equation}
Then for any generator $G=\alpha L+\beta B_X+\gamma B_Y\in\mathfrak{so}(2,1)$ (possibly time-dependent, as in the perturbed flows of \cite{Ebert2026Cone}), and $\hat G:=\alpha i/2+\beta j/2+\gamma k/2$,
\begin{equation}
\dot V=[\hat G,V] \iff \begin{pmatrix}\dot X\\ \dot Y\\ \dot T\end{pmatrix}=G\begin{pmatrix}X\\Y\\T\end{pmatrix}. \tag{20}
\end{equation}
Moreover $V\bar V=T^2-X^2-Y^2\equiv0$ on the cone: the cone identity is exactly the statement that $V$ is a \emph{null} split quaternion, consistent with the rank-one spinor picture of \cite{Ebert2026EigenCoord}.
\end{theorem}

\begin{proof}
Both claims are direct computations in $\mathbb H'$: expand $\hat GV-V\hat G$ in the basis $1,i,j,k$ and compare coefficients to $G(X,Y,T)^\top$ using the explicit matrices $L,B_X,B_Y$ of \cite{Ebert2026EigenCoord}; expand $V\bar V=(Ti+Yj-Xk)(-Ti-Yj+Xk)$ using $i^2=-1$, $j^2=k^2=1$, and the vanishing of all cross terms (the three imaginary units pairwise anticommute, so every cross term cancels against its mirror). We verified both symbolically for general $\alpha,\beta,\gamma$, with zero residual. The particular sign pattern in (19) --- $Y$ on $j$ but $-X$ on $k$ --- is forced: it is the unique (up to overall sign) assignment of $(X,Y,T)$ to $(i,j,k)$ for which (20) holds; the naive guess $V=Ti+Xj+Yk$ does not satisfy it.
\end{proof}

\begin{remark}[Why the clock does not join this equation]
One might hope to extend (19)--(20) to the full state by appending the clock as a scalar, $Q=\Theta+V$ for $\Theta\in\{\nu,E,M\}$, and ask whether $\dot Q=[\hat G,Q]$ reproduces $\dot\Theta$ alongside the shape flow. It cannot: the commutator of any two split quaternions has identically zero scalar part, since every basis commutator $[i,j],[j,k],[k,i]$ is itself a pure imaginary unit. Hence $[\hat G,Q]=[\hat G,V]$ for \emph{every} choice of $\Theta$, regardless of its own dynamics --- the commutator structure that reproduces the shape flow is, by construction, blind to any scalar coordinate riding along with it. A one-sided product $\dot Q=Q\Omega$ can be made to reproduce a prescribed $\dot\Theta$, but only by letting $\Omega$ depend on the state however is needed to match (15) or $\dot M=\sqrt{\mu/T^3}$; this repackages Theorem 3.1 rather than extending it, and is not offered as a result here. The obstruction is structural, not a failure of search: only the shape $(X,Y,T)$ carries the $\mathfrak{so}(2,1)$ symmetry of \cite{Ebert2026Cone,Ebert2026EigenCoord}, and (19)--(20) show that symmetry's natural algebraic home is the \emph{pure} part of $\mathbb H'$, to which the scalar part is, by the algebra's own multiplication table, blind. This is an algebraic proof of the same fact Section 1 states dynamically: the clock is external to the cone's own symmetry, not a coordinate that symmetry secretly already carries.
\end{remark}

\section{Discussion}

The cone construction of \cite{Ebert2026Cone} isolates orbital shape --- energy, angular momentum, mass, eccentricity --- as geometric directions on $X^2+Y^2=T^2$, but by construction says nothing about time-of-flight along a given orbit. This note closes that gap with two short, exact steps: rewriting the orbit equation in cone coordinates (Proposition 2.1), and combining it with the cone's own angular-momentum formula to obtain a closed-form clock equation (Theorem 3.1). The result is a genuine extension of the construction's reach --- from a static description of shape to a coupled description of shape \emph{and} position-in-time --- rather than a reinterpretation of any coordinate already carrying a physical meaning in \cite{Ebert2026Cone}. In particular, $T$ remains the semi-major axis throughout; time enters only as the independent variable of the new ODE (15), exactly as it already does in the shape-evolution equations of \cite{Ebert2026Cone}.

Sections 5 and 6 push the clock further into the algebraic machinery of \cite{Ebert2026Table,Ebert2026EigenCoord} and find a sharp two-part answer. On one hand, the classical $E$--$\nu$ correspondence is not merely algebraically convenient but \emph{is}, exactly, a hyperbolic Möbius transformation generated by the cone's own $B_X$, at rapidity $\tfrac12\operatorname{artanh}(e)$ (Proposition 5.1) --- a genuinely new addition to the clock, not a restatement of Theorem 3.1, and one more instance of the half-angle pattern already forced elsewhere in this project. On the other hand, the full shape-evolution system $(\dot X,\dot Y,\dot T)$ collapses to a single split-quaternion commutator equation (Theorem 6.1), a literal Lax-pair form for the ``Lie-algebra flow'' language of Section 1, while the clock variable is provably excluded from that same equation by an elementary fact about quaternion commutators (Remark following Theorem 6.1). Read together, these results say that the clock is connected to the cone's symmetry algebra through its \emph{kinematics} ($E\leftrightarrow\nu$, Proposition 5.1) but not through its \emph{dynamics} (the shape flow, Theorem 6.1) --- a distinction this note could not have stated precisely before working out both directions.

\section*{AI Assistance Disclosure}

Large language model tools (Anthropic's Claude) were used in the derivation, cross-checking, and LaTeX preparation of this note, including symbolic verification of Proposition 2.1 and numerical verification of Theorem 3.1 against independent Cartesian integration. Sections 5 and 6 were developed collaboratively with Claude: the split-quaternion multiplication table, the generator dictionary (18), the commutator identity (20), the correct sign assignment in (19), and the $E$--$\nu$ Möbius identity (16)--(17) were all checked by direct symbolic computation (computer algebra), and Proposition 5.1 was additionally verified numerically to $30$-digit precision against the orbit and epoch of Section 4. All mathematical claims were independently verified by the author.

\begin{thebibliography}{9}
\bibitem{Ebert2026Cone} J.~Ebert, \emph{The Dynamic Slicing Operator on the AM--GM Cone: A Quadric Geometric Framework for Keplerian Orbits and Perturbed Trajectories}, Preprint (2026).
\bibitem{Ebert2026Table} J.~Ebert, \emph{The Cutting Plane: Every Line Through the Multiplication Table Is a Conic Section}, Preprint (2026).
\bibitem{Ebert2026EigenCoord} J.~Ebert, \emph{Every Cut Is an Orbit: Eigen-Coordinates and a Power-Map Companion to the Cutting-Plane Theorem}, Preprint (2026).
\end{thebibliography}

\end{document}
